\documentclass[12pt, a4paper]{article}

% ============================================================
%  Paquetes
% ============================================================
\usepackage[utf8]{inputenc}
\usepackage[T1]{fontenc}
\usepackage[spanish]{babel}
\usepackage{amsmath, amssymb}
\usepackage[ruled, vlined, linesnumbered, spanish]{algorithm2e}
\usepackage{geometry}
\usepackage{hyperref}

\geometry{margin=2.5cm}

% ============================================================
%  Metadatos del documento
% ============================================================
\title{Búsqueda Tabú Simple y Reactiva --- Pseudocódigo}
\author{}
\date{}

% ============================================================
\begin{document}
\maketitle

% ------------------------------------------------------------
\section{Notación común}
% ------------------------------------------------------------

A continuación se listan los símbolos compartidos por la Búsqueda Tabú Simple
(TS) y la Búsqueda Tabú Reactiva (RTS). Estos símbolos coinciden con los del
documento de pseudocódigo de Recocido Simulado y de Abejas Artificiales del
proyecto MetaCARP.

\begin{itemize}
    \item $G = (V, E, R)$: grafo de la instancia CARP. $V$ son los nodos, $E$ las
          aristas y $R \subseteq E$ las aristas requeridas.
    \item $s$: solución --- conjunto de rutas $\{r_1, r_2, \ldots, r_k\}$, cada
          ruta es una secuencia de etiquetas de tareas requeridas con el depósito
          al inicio y al final.
    \item $s^{*}$: mejor solución encontrada durante toda la corrida.
    \item $c(s)$: costo puro (suma de costos de travesía).
    \item $\mathrm{viol}(s)$: violación total de capacidad (exceso de demanda).
    \item $f(s) = c(s) + \lambda \cdot \mathrm{viol}(s)$: objetivo penalizado
          ($\lambda > 0$).
    \item $\mathcal{N}(s)$: vecindario de $s$ inducido por los operadores
          activos.
    \item $\mathcal{O}_{\mathrm{intra}} = \{\texttt{relocate\_intra},\,
          \texttt{swap\_intra},\, \texttt{2opt\_intra}\}$: operadores
          intra-ruta.
    \item $\mathcal{O}_{\mathrm{inter}} = \{\texttt{relocate\_inter},\,
          \texttt{swap\_inter},\, \texttt{2opt\_star},\,
          \texttt{cross\_exchange},\, \texttt{or\_opt\_2},\,
          \texttt{or\_opt\_3}\}$: operadores inter-ruta.
    \item $\mathcal{T}$: lista tabú (cola FIFO de claves de movimientos).
    \item $\tau$: tenencia tabú (capacidad de $\mathcal{T}$). En TS es fija; en
          RTS es dinámica y vive en $[\tau_{\min}, \tau_{\max}]$.
    \item $n = |R|$: número de tareas requeridas (variable de escala
          instance-aware).
    \item $p_{\mathrm{inter}} \in [0,1]$: probabilidad de elegir el grupo
          inter-ruta cuando la solución actual es factible.
    \item $\alpha_{\mathrm{inter}} \in [0,1]$: probabilidad de elegir el grupo
          inter-ruta cuando la solución actual viola capacidad.
    \item $H(s)$: hash canónico de una solución (en RTS, para detectar ciclos).
\end{itemize}

\medskip
\noindent
\textbf{Clave tabú $\kappa(m)$.} El movimiento $m$ aplicado para pasar de $s$
a $s'$ se identifica por la tupla
$\kappa(m) = (\mathrm{operador},\, r_a, r_b,\, i, j, k, l,\, \text{etiquetas
movidas})$. Dos movimientos con la misma clave se consideran iguales.

\medskip
\noindent
\textbf{Sesgo inter/intra (compartido con SA y ABC).} En cada iteración del
bucle principal se calcula $v \leftarrow \mathrm{viol}(s)$ y se define
\[
p_{\mathrm{efec}} \;=\;
\begin{cases}
  \alpha_{\mathrm{inter}} & \text{si } v > 0, \\
  p_{\mathrm{inter}}     & \text{si } v = 0.
\end{cases}
\]
Con probabilidad $p_{\mathrm{efec}}$ se elige el grupo
$\mathcal{O}_{\mathrm{inter}}$; en caso contrario,
$\mathcal{O}_{\mathrm{intra}}$. Dentro del grupo elegido, el operador concreto
se muestrea de manera uniforme.

% ------------------------------------------------------------
\section{Búsqueda Tabú Simple (TS)}
% ------------------------------------------------------------

Implementación más simple del Tabu Search clásico de Glover (1986):
best-improvement sobre un lote muestreado del vecindario, lista tabú FIFO de
tenencia fija, criterio de aspiración clásico y doble criterio de parada
(iteraciones máximas o estancamiento). La lista tabú se mantiene mediante una
deque acompañada de un \emph{set} y de un \emph{Counter} paralelo para que las
consultas y los descartes por overflow se realicen en $O(1)$.

\begin{algorithm}[H]
\DontPrintSemicolon
\caption{Búsqueda Tabú Simple para CARP}
\label{alg:ts_simple}

\KwIn{Instancia CARP $G$;\; solución inicial $s_0$;\;
      $\mathit{iteraciones\_max}$;\; $\mathit{max\_iter\_sin\_mejora}$;\;
      $\mathit{tam\_vecindario}$;\; tenencia tabú $\tau$;\;
      $p_{\mathrm{inter}}$,\; $\alpha_{\mathrm{inter}}$;\;
      operadores $\mathcal{O}$}
\KwOut{Mejor solución encontrada $s^{*}$}

\BlankLine
\tcp{Inicialización}
$s \leftarrow s_0$;\quad $s^{*} \leftarrow s$;\quad
$c^{*} \leftarrow c(s)$ \;
$\mathcal{T} \leftarrow$ deque FIFO de capacidad $\tau$ (acompañada de un set
y de un Counter paralelos para lookup y descarte en $O(1)$) \;
$\mathit{iter} \leftarrow 0$;\quad $\mathit{iter\_sin\_mejora} \leftarrow 0$ \;
$\mathit{aspiraciones} \leftarrow 0$;\quad
$\mathit{todos\_tabu} \leftarrow 0$ \;

\BlankLine
\While{$\mathit{iter} < \mathit{iteraciones\_max}$}{
    \If{$\mathit{iter\_sin\_mejora} \geq \mathit{max\_iter\_sin\_mejora}$}{
        \textbf{break}
            \tcp{parada anticipada por estancamiento}
    }

    \BlankLine
    \tcp{Paso 1: selección del grupo de operadores (sesgo inter/intra)}
    $v \leftarrow \mathrm{viol}(s)$ \;
    \eIf{$v > 0$}{
        $p_{\mathrm{efec}} \leftarrow \alpha_{\mathrm{inter}}$
    }{
        $p_{\mathrm{efec}} \leftarrow p_{\mathrm{inter}}$
    }
    \eIf{$\mathrm{Uniforme}(0,1) < p_{\mathrm{efec}}$}{
        $\mathcal{O}_{\mathrm{grupo}} \leftarrow \mathcal{O}_{\mathrm{inter}}$
    }{
        $\mathcal{O}_{\mathrm{grupo}} \leftarrow \mathcal{O}_{\mathrm{intra}}$
    }

    \BlankLine
    \tcp{Paso 2: generación del lote de vecinos}
    $\mathcal{N} \leftarrow \emptyset$;\quad
    $\mathcal{M} \leftarrow \emptyset$
        \tcp{vecinos y movimientos asociados}
    \For{$k \leftarrow 1$ \KwTo $\mathit{tam\_vecindario}$}{
        $\mathit{op} \leftarrow$ \textsc{SelectUniform}($\mathcal{O}_{\mathrm{grupo}}$) \;
        $(s'_k, m_k) \leftarrow$ \textsc{AplicarOperador}($\mathit{op},\, s$) \;
        $\mathcal{N} \leftarrow \mathcal{N} \cup \{s'_k\}$;\quad
        $\mathcal{M} \leftarrow \mathcal{M} \cup \{m_k\}$ \;
    }

    \BlankLine
    \tcp{Paso 3: evaluación vectorizada del lote}
    Evaluar $f(s'_k)$ para todo $k$ en una sola pasada NumPy
        (\texttt{costo\_lote\_ids}) \;

    \BlankLine
    \tcp{Paso 4: mejor vecino admisible (no-tabú o tabú con aspiración)}
    $\mathit{ind\_admisible} \leftarrow -1$;\quad
    $\mathit{ind\_total} \leftarrow 0$ \;
    $\mathit{aspiracion\_iter} \leftarrow \textsc{Falso}$ \;
    \ForEach{$k \in \{1, \ldots, \mathit{tam\_vecindario}\}$}{
        Actualizar el mejor total $\mathit{ind\_total}$ con $\arg\min f(s'_k)$ \;
        $\mathit{es\_tabu} \leftarrow (\kappa(m_k) \in \mathcal{T})$ \;
        $\mathit{aspiracion} \leftarrow (f(s'_k) < c^{*} - \varepsilon)$
            \tcp{aspiración clásica}
        \If{$\mathit{es\_tabu}$ \textbf{ y } \textbf{no} $\mathit{aspiracion}$}{
            \textbf{continue}
                \tcp{movimiento tabú sin aspiración: se descarta}
        }
        Actualizar $\mathit{ind\_admisible}$ con $\arg\min f(s'_k)$ \;
        \If{$\mathit{es\_tabu}$ \textbf{ y } $\mathit{aspiracion}$}{
            $\mathit{aspiracion\_iter} \leftarrow \textsc{Verdadero}$ \;
        }
    }

    \BlankLine
    \tcp{Paso 5: elección final del movimiento}
    \eIf{$\mathit{ind\_admisible} = -1$}{
        $k^{*} \leftarrow \mathit{ind\_total}$;\quad
        $\mathit{todos\_tabu} \leftarrow \mathit{todos\_tabu} + 1$
            \tcp{todos tabú: fallback al mejor tabú}
    }{
        $k^{*} \leftarrow \mathit{ind\_admisible}$ \;
        \lIf{$\mathit{aspiracion\_iter}$}{
            $\mathit{aspiraciones} \leftarrow \mathit{aspiraciones} + 1$
        }
    }
    $s \leftarrow s'_{k^{*}}$;\quad $m^{*} \leftarrow m_{k^{*}}$
        \tcp{TS avanza SIEMPRE, aunque empeore}

    \BlankLine
    \tcp{Paso 6: actualización FIFO de la lista tabú con Counter paralelo}
    \If{$|\mathcal{T}| = \tau$}{
        $\kappa_{\mathrm{desc}} \leftarrow$ clave más antigua de $\mathcal{T}$ \;
        $\mathrm{conteo}[\kappa_{\mathrm{desc}}] \leftarrow
            \mathrm{conteo}[\kappa_{\mathrm{desc}}] - 1$ \;
        \If{$\mathrm{conteo}[\kappa_{\mathrm{desc}}] = 0$}{
            eliminar $\kappa_{\mathrm{desc}}$ del set y del Counter
                \tcp{descarte $O(1)$}
        }
    }
    Insertar $\kappa(m^{*})$ en la deque, en el set, e incrementar
        $\mathrm{conteo}[\kappa(m^{*})]$ \;

    \BlankLine
    \tcp{Paso 7: actualización del mejor global}
    \eIf{$f(s) < c^{*} - \varepsilon$}{
        $s^{*} \leftarrow s$;\quad $c^{*} \leftarrow f(s)$;\quad
        $\mathit{iter\_sin\_mejora} \leftarrow 0$ \;
    }{
        $\mathit{iter\_sin\_mejora} \leftarrow \mathit{iter\_sin\_mejora} + 1$ \;
    }
    $\mathit{iter} \leftarrow \mathit{iter} + 1$ \;
}

\BlankLine
\Return $s^{*}$
\end{algorithm}

\bigskip
\noindent
\textbf{Parámetros por defecto (TS simple).} La implementación canónica del
proyecto utiliza $\mathit{iteraciones\_max} = 400$,
$\mathit{max\_iter\_sin\_mejora} = 100$, $\mathit{tam\_vecindario} = 25$,
$\tau = 20$. Para el sesgo: $\alpha_{\mathrm{inter}} = 0.8$ y
$p_{\mathrm{inter}} = 0.6$ (los mismos valores que SA y RTS, para garantizar
comparabilidad entre las tres metaheurísticas).

% ------------------------------------------------------------
\section{Búsqueda Tabú Reactiva (RTS)}
% ------------------------------------------------------------

Reactive Tabu Search de Battiti \& Tecchiolli (1994). Extiende TS simple con
tres mecanismos adaptativos: (i) tenencia tabú dinámica, (ii) memoria de
soluciones visitadas por hash canónico para detectar ciclos, y
(iii) mecanismo de escape (diversificación fuerte) cuando una solución se
visita demasiadas veces. Hereda de TS simple la lista tabú FIFO, el
best-improvement sobre lote muestreado, la aspiración clásica y el doble
criterio de parada.

\medskip
\noindent
\textbf{Hash canónico $H(s)$.} Se ordenan las rutas lexicográficamente como
tuplas (los vehículos son intercambiables) y se conserva el orden interno de
cada ruta (CARP puede ser dirigido). El hash se calcula sobre la tupla de
tuplas resultante. Esto permite detectar ciclos en $O(1)$ vía diccionario.

\medskip
\noindent
\textbf{Parámetros instance-aware (cálculo automático cuando el usuario no
los especifica).} Sea $n = |R|$. La implementación calcula:
\[
\begin{aligned}
  \mathit{iteraciones\_max}     &\leftarrow \max(50,\; 20 n) \\
  \mathit{max\_iter\_sin\_mejora} &\leftarrow \max(20,\; 5 n) \\
  \mathit{tam\_vecindario}      &\leftarrow \max(20,\; 2 n) \\
  \tau_{\mathrm{ini}}            &\leftarrow \max(3,\; \mathrm{round}(\sqrt{n})) \\
  \tau_{\min}                    &\leftarrow 3 \\
  \tau_{\max}                    &\leftarrow \max(15,\; \mathrm{round}(3\sqrt{n})) \\
  \mathit{iter\_paciencia}       &\leftarrow \max(5,\; \mathrm{round}(2\sqrt{n})) \\
  \mathit{num\_mov\_escape}      &\leftarrow \max(3,\; n/10).
\end{aligned}
\]
El piso $\tau_{\max} \geq 15$ garantiza un rango útil incluso para instancias
muy pequeñas (p.\,ej.\ \texttt{gdb1} con $n \approx 11$). Adicionalmente,
$\mathit{factor\_aumento} = 1.2$, $\mathit{factor\_reduccion} = 0.9$ y
$\mathit{umbral\_repeticiones\_escape} = 3$.

\begin{algorithm}[H]
\DontPrintSemicolon
\caption{Búsqueda Tabú Reactiva (RTS) para CARP}
\label{alg:rts}

\KwIn{Instancia CARP $G$;\; solución inicial $s_0$;\;
      $\mathit{iteraciones\_max}$,\; $\mathit{max\_iter\_sin\_mejora}$,\;
      $\mathit{tam\_vecindario}$;\;
      $\tau_{\mathrm{ini}}, \tau_{\min}, \tau_{\max}$;\;
      $\mathit{factor\_aumento} > 1$,\;
      $0 < \mathit{factor\_reduccion} < 1$;\;
      $\mathit{iter\_paciencia}$;\;
      $\mathit{umbral\_repeticiones\_escape}$,\;
      $\mathit{num\_mov\_escape}$;\;
      $p_{\mathrm{inter}}$,\; $\alpha_{\mathrm{inter}}$;\;
      operadores $\mathcal{O}$}
\KwOut{Mejor solución encontrada $s^{*}$}

\BlankLine
\tcp{Inicialización}
$s \leftarrow s_0$;\quad $s^{*} \leftarrow s$;\quad $c^{*} \leftarrow c(s)$ \;
$\tau \leftarrow \tau_{\mathrm{ini}}$ \;
$\mathcal{T} \leftarrow$ deque FIFO de capacidad $\tau$ (con set y Counter paralelos) \;
$\mathit{historial} \leftarrow \{H(s) \mapsto (\text{última vista}=-1,\,
    \text{veces vistas}=1)\}$ \;
$\mathit{iter\_sin\_repeticion} \leftarrow 0$;\quad
$\mathit{iter\_sin\_mejora} \leftarrow 0$ \;
$\mathit{iter} \leftarrow 0$ \;

\BlankLine
\While{$\mathit{iter} < \mathit{iteraciones\_max}$}{
    \lIf{$\mathit{iter\_sin\_mejora} \geq \mathit{max\_iter\_sin\_mejora}$}{
        \textbf{break}
    }

    \BlankLine
    \tcp{Pasos 1--5: idénticos a TS simple (sesgo, lote, evaluación,
        best-non-tabu con aspiración, fallback "todos tabú")}
    Ejecutar los pasos 1--5 del Algoritmo~\ref{alg:ts_simple}
        sobre $s$, $\mathcal{T}$ y $\mathcal{O}$ \;
    $s \leftarrow s'_{k^{*}}$;\quad $m^{*} \leftarrow m_{k^{*}}$ \;

    \BlankLine
    \tcp{Paso 6: insertar el movimiento en la lista tabú FIFO}
    Insertar $\kappa(m^{*})$ en $\mathcal{T}$ con el mismo protocolo de
        descarte del Algoritmo~\ref{alg:ts_simple}
        (deque + set + Counter paralelos) \;

    \BlankLine
    \tcp{Paso 7: actualización del mejor global}
    \If{$f(s) < c^{*} - \varepsilon$}{
        $s^{*} \leftarrow s$;\quad $c^{*} \leftarrow f(s)$;\quad
        $\mathit{iter\_sin\_mejora} \leftarrow 0$ \;
    }
    \lElse{
        $\mathit{iter\_sin\_mejora} \leftarrow \mathit{iter\_sin\_mejora} + 1$
    }

    \BlankLine
    \tcp{Paso 8: MECANISMO REACTIVO DEL TENURE}
    $h \leftarrow H(s)$ \;
    \eIf{$h \in \mathit{historial}$}{
        \tcp{repetición detectada: posible ciclo}
        $\mathit{historial}[h].\text{veces} \leftarrow
            \mathit{historial}[h].\text{veces} + 1$ \;
        $\mathit{iter\_sin\_repeticion} \leftarrow 0$ \;

        \BlankLine
        \tcp{Aumento del tenure (siempre con +1 mínimo si hay margen)}
        $\tau_{\mathrm{nuevo}} \leftarrow
            \min\!\big(\tau_{\max},\, \max(\tau + 1,\, \mathrm{round}(\tau
            \cdot \mathit{factor\_aumento}))\big)$ \;
        \If{$\tau_{\mathrm{nuevo}} > \tau$}{
            $\tau \leftarrow \tau_{\mathrm{nuevo}}$;\quad
            reconstruir $\mathcal{T}$ con la nueva capacidad
                \tcp{\texttt{deque.maxlen} es solo lectura}
        }

        \BlankLine
        \tcp{ESCAPE si la solución se repitió suficientes veces}
        \If{$\mathit{historial}[h].\text{veces} \geq
            \mathit{umbral\_repeticiones\_escape}$}{
            \For{$j \leftarrow 1$ \KwTo $\mathit{num\_mov\_escape}$}{
                $\mathit{op} \leftarrow$ \textsc{SelectUniform}($\mathcal{O}$)
                    \tcp{lista completa, sin sesgo durante el escape} \;
                $s \leftarrow$ \textsc{AplicarOperador}($\mathit{op},\, s$)
                    \tcp{lista tabú ignorada durante el escape}
            }
            $\mathcal{T} \leftarrow \emptyset$;\quad
                $\mathit{historial} \leftarrow \emptyset$
                \tcp{limpieza tras escape (memoria irrelevante en la zona nueva)} \;
            $\mathit{historial}[H(s)] \leftarrow (\text{última vista}=\mathit{iter},\,
                \text{veces vistas}=1)$ \;
            $\mathit{iter\_sin\_repeticion} \leftarrow 0$ \;
            \lIf{$f(s) < c^{*} - \varepsilon$}{
                $s^{*} \leftarrow s$;\quad $c^{*} \leftarrow f(s)$
            }
            $\mathit{iter\_sin\_mejora} \leftarrow 0$
                \tcp{reset INCONDICIONAL: el escape merece presupuesto completo} \;
        }
    }{
        \tcp{solución nueva: registrar e incrementar paciencia}
        $\mathit{historial}[h] \leftarrow (\text{última vista}=\mathit{iter},\,
            \text{veces vistas}=1)$ \;
        $\mathit{iter\_sin\_repeticion} \leftarrow
            \mathit{iter\_sin\_repeticion} + 1$ \;
        \BlankLine
        \tcp{Reducción del tenure tras suficientes iteraciones sin repeticiones}
        \If{$\mathit{iter\_sin\_repeticion} \geq \mathit{iter\_paciencia}$}{
            $\tau_{\mathrm{nuevo}} \leftarrow
                \max\!\big(\tau_{\min},\, \min(\tau - 1,\, \mathrm{round}(\tau
                \cdot \mathit{factor\_reduccion}))\big)$ \;
            \If{$\tau_{\mathrm{nuevo}} < \tau$}{
                $\tau \leftarrow \tau_{\mathrm{nuevo}}$;\quad
                reconstruir $\mathcal{T}$ truncando los elementos más antiguos \;
            }
            $\mathit{iter\_sin\_repeticion} \leftarrow 0$
                \tcp{reinicio incondicional: evita acumular paciencia} \;
        }
    }

    \BlankLine
    $\mathit{iter} \leftarrow \mathit{iter} + 1$ \;
}

\BlankLine
\Return $s^{*}$
\end{algorithm}

\bigskip
\noindent
\textbf{Decisiones de diseño.}
\begin{itemize}
    \item \emph{Hash canónico:} ordenar rutas lexicográficamente (los vehículos
          son intercambiables) sin invertir el orden interno (CARP puede ser
          dirigido). Falsos positivos descartados por construcción.
    \item \emph{Aumento del tenure:} $\tau_{\mathrm{nuevo}} =
          \min(\tau_{\max},\, \max(\tau + 1,\,
          \mathrm{round}(\tau \cdot \mathit{factor\_aumento})))$. El
          $\max(\tau+1, \cdot)$ garantiza al menos +1 cuando hay margen,
          incluso si $\mathrm{round}(\tau \cdot 1.05) = \tau$ en instancias
          pequeñas.
    \item \emph{Reducción del tenure:} simétrica, $\tau_{\mathrm{nuevo}} =
          \max(\tau_{\min},\, \min(\tau - 1,\,
          \mathrm{round}(\tau \cdot \mathit{factor\_reduccion})))$ tras
          $\mathit{iter\_paciencia}$ iteraciones sin repeticiones; el
          contador se reinicia incondicionalmente para evitar acumular
          paciencia indefinidamente.
    \item \emph{Escape:} aplica $\mathit{num\_mov\_escape}$ movimientos
          aleatorios consecutivos con la lista completa de operadores
          ignorando $\mathcal{T}$; luego limpia $\mathcal{T}$ y
          $\mathit{historial}$ (la memoria de la zona anterior ya no es
          relevante en la zona nueva) y resetea $\mathit{iter\_sin\_mejora}$
          incondicionalmente para dar a la nueva región un presupuesto
          completo de explotación.
    \item \emph{Reescalado de $\mathcal{T}$:} como \texttt{deque.maxlen} es de
          solo lectura, cada cambio de tenure reconstruye la deque
          preservando los elementos más recientes y manteniendo el set y el
          Counter paralelos en sincronía.
\end{itemize}

\end{document}
